% SGA 3，第 VII_B 篇，英文重建稿——第 1 轮。
% 范围：权威底本分篇第 90--91 页，第 5.2.1.A 节全文，止于第 5.2.1.B 节之前。
% 合并阅读本页码： 603--604.
% 再版印刷页码： 592--593.
% 本范围内编者注： 155.
% 权威底本： Polo--Gille Exp7B-13oct24.pdf,
% SHA-256 D303690E0195654BBE7C3DAA5D89EA23E24CD709A3575217234AC19AE8836694.
% 曾参考 Jacob Reinhold 的 CC BY 4.0 Markdown 作为比对谱系。
% 原英文 Loop 1 阶段说明：两幅图采用直接取自 PDF 的精确 PNG 裁片。当前中文工程已采用原生图。

\paragraph{\SGARefLink{sga3:VIIB:5.2.1.A}{5.2.1.A}.}
\phantomsection
\label{sga3:VIIB:5.2.1.A}
\footnote{\label{sga3:VIIB:nde:155}编者增加了 \SGARefLink{sga3:VIIB:5.2.1.A}{5.2.1.A} 和 \SGARefLink{sga3:VIIB:5.2.1.B}{5.2.1.B}。}
对每个 \(r\in\mathbb N\)，置
\[
  \mathcal I_r=\overline{I^2}+I_r,
\]
于是
\[
  \mathcal I_r/\overline{I^2}\simeq\omega_r.
\]
此时有如下各行正合的交换图：
\begin{center}
  \phantomsection
  \label{sga3:VIIB:5.2.1.A:exact-rows}
  % SGA3-VII-DIAG-130
  % SGA3-VII-DIAG-130
% Authority: Exp7B-13oct24.pdf, local p. 90.
% Crop: figures/Exp7B_localp090_521A_exact_rows.png
% Locator: tex/components/84_expose_VIIB_section521A_pseudobasis_splittings_en.tex, Loop 1 line 28.
\begin{tikzcd}[row sep=3em,column sep=2.6em]
0 \arrow[r] & \mathcal I_{r-1} \arrow[r] \arrow[d]
  & \mathcal I_r \arrow[r] \arrow[d]
  & \mathcal I_r/\mathcal I_{r-1}
      \arrow[r] \arrow[d,"\wr"]
  & 0 \\
0 \arrow[r] & \omega_{r-1} \arrow[r]
  & \omega_r \arrow[r]
  & \omega_r/\omega_{r-1} \arrow[r]
  & 0
\end{tikzcd}

\end{center}

由于 \(k\) 为域，各行分裂。因此，\(\mathcal I_{r-1}\) 的伪基 \(\mathscr B_{r-1}\) 可补成 \(\mathcal I_r\) 的伪基
\[
  \mathscr B_{r-1}\cup\mathscr B'_r
\]
于是，以 \(\mathscr B'_r\) 为伪基的闭子空间 \(S_r\) 为 \(\mathcal I_{r-1}\) 在 \(\mathcal I_r\) 中的补空间，而投影
\[
  \pi:I\longrightarrow\omega
\]
诱导 \(S_r\) 到 \(\omega_{r-1}\) 在 \(\omega_r\) 中的一个补空间 \(\omega'_r\) 的同构。

令 \(\mathcal I_\infty\) 为闭理想
\[
  \mathcal I_\infty
  =
  \overline{\bigcup_r\mathcal I_r}.
\]
它同样在 \(I\) 中有一个补空间 \(S_\infty\)，而 \(\pi\) 诱导
\[
  \mathcal I_\infty/\overline{I^2}
\]
到各 \(\omega_r\) 的并的闭包
\[
  \omega_\infty
  =
  \overline{\bigcup_r\omega_r}
\]
的同构，以及 \(S_\infty\) 到 \(\omega_\infty\) 在 \(\omega\) 中的一个补空间 \(\omega'_\infty\) 的同构。将后一个同构记为
\[
  \eta:S_\infty\xrightarrow{\sim}\omega'_\infty.
\]
于是得到连续线性映射：
\begin{center}
  \phantomsection
  \label{sga3:VIIB:5.2.1.A:splitting-maps}
  % SGA3-VII-DIAG-131
  % SGA3-VII-DIAG-131
% Authority: Exp7B-13oct24.pdf, local p. 90.
% Crop: figures/Exp7B_localp090_521A_splitting_maps.png
% Locator: tex/components/84_expose_VIIB_section521A_pseudobasis_splittings_en.tex, Loop 1 line 74.
\begin{tikzcd}[row sep=3.2em,column sep=3em]
\overline{I^2}\times S_\infty\times\widehat{\bigoplus}_r S_r
  \arrow[r,"\phi"]
  \arrow[d,"\eta" left]
  \arrow[d,phantom,"{\times\theta}" right]
  & I \\
\omega=\omega'_\infty\times\omega_\infty & {}
\end{tikzcd}

\end{center}

这里
\[
  \widehat{\bigoplus}_r S_r
\]
为 \(\operatorname{PC}(k)\) 中各 \(S_r\) 的直和，即
\[
  \left(\prod_r S_r^\dagger\right)^*
\]
\textup{（参见 \SGARefLink{sga3:VIIB:0.2.2}{0.2.2} 的编者注 \textup{(16)}）}，且
\[
  \theta:
  \widehat{\bigoplus}_r S_r
  \longrightarrow\omega_\infty
\]
由映射
\[
  S_r\xrightarrow{\sim}\omega'_r
  \lhook\joinrel\longrightarrow\omega_\infty.
\]
诱导。由此可见，得到所需类型的截面
\[
  \sigma:\omega\longrightarrow I
\]
的一个充分条件\textup{（但并非必要条件，见下文）}是 \(\theta\) 为同构。由对偶性\textup{（参见 \SGARefLink{sga3:VIIB:0.2.2}{0.2.2}）}，这等价于线性映射
\[
  \omega_\infty^\dagger
  \longrightarrow
  \prod_r S_r^\dagger
\]
为双射。
